\documentclass[11pt,a4paper]{article}
\usepackage[T1]{fontenc}
\usepackage[utf8]{inputenc}
\usepackage{lmodern}
\usepackage[margin=25mm]{geometry}
\usepackage{amsmath,amssymb,booktabs,tabularx,xurl,hyperref}
\hypersetup{colorlinks=true,linkcolor=blue!50!black,urlcolor=blue!50!black,
  pdftitle={Keyed Improvement: Source Updates, Path Semantics and Lattice Closure},
  pdfauthor={Mettapedia Working Notes}}
\setlength{\emergencystretch}{3em}
\setlength{\parindent}{0pt}
\setlength{\parskip}{5pt}
\newcommand{\code}[1]{\nolinkurl{#1}}
\title{Keyed Improvement\\
  \large Source Updates, Path Semantics and Lattice Closure}
\author{Mettapedia Working Notes}
\date{31 October 2026}
\begin{document}
\maketitle

\begin{abstract}
Seven Lean modules connect a source-shaped remove/add improvement to keyed
minimum replacement, and an improvement-driven shortest-path worklist to an
independent weighted-path semantics. The worklist terminates over a finite
vertex universe with nonnegative natural costs. Its final labels equal the
least closure in the reverse numerical order, using the existing chaotic
iteration theory. Source-shaped adjacency loops admit a stronger trace law:
their once-read source distance can be retained across their own updates.
Final-distance confluence is kept separate from propagation counts, source
round counters, concurrent publication and native implementation costs.
\end{abstract}

\section{Physical deletion can denote monotone information improvement}

The inspected shortest-path program keeps at most one distance row per key.
Its \code{improve} operation inserts an absent distance, or removes the old
row and inserts a strictly smaller distance. A successful improvement returns
the key; otherwise it returns no answer. Its surrounding driver builds a new
frontier and ultimately returns a round counter as well as exposing distances.

\code{KeyedImprovement} models a slot as \code{Option Value} and independently
models its finite row relation. Encoding a slot store yields one row per
occupied key. The proved row law preserves the entire encoded relation and
the strict-improvement flag after a completed update. It requires the
functional-relation invariant and the correct observed old value.

For exact natural distances, absence is genuine $\infty$. The keyed update is
\[
  d[k]\gets\min(d[k],c).
\]
It is the cap operation in the existing \code{TropicalAffineSummary} algebra.
Thus removing and replacing a physical row can represent a monotone semantic
operation. Monotonicity belongs to the information order, not to the literal
set of physical rows.

Fixed candidate updates commute on final values. Their changed/not-changed
flags do not: starting at $10$, applying $7$ then $5$ changes twice, while
applying $5$ then $7$ changes once. Edge-derived candidates also depend on
other live labels; their individual transitions are not asserted to commute.

\section{Publication and captured observations}

A raw remove/add sequence has an observable intermediate absence. Two
threads using the same stale old row can additionally destroy the one-value
invariant: after the row $(k,10)$ becomes $(k,7)$, erasing the nonexistent
$(k,10)$ and adding $(k,5)$ leaves both $(k,7)$ and $(k,5)$.

\code{KeyedImprovementTransactions} prepares a key, candidate and observed
value. Its logical publication validates the observed value against the
current key before applying the prepared branch. A stale ticket returns
retry, not ``no improvement.'' A fallback recomputes against the live store,
and is proved equal to ordinary live improvement. Unrelated key updates need
not force retry.

These are value-level publication laws. A native implementation must supply
the corresponding atomicity, failure handling and ownership. If occurrence
identity is observable, validation must also retain the required occurrence
or generation identity; equality of payload values does not certify pointer
lifetime. Existing captured-reader laws remain relevant to readers surviving
the physical row replacement.

\section{Independent shortest-path meaning}

An arc has a source, target and natural weight. Parallel edges and zero
weights are permitted. \code{WeightedPathRelaxation} reuses Mathlib's quiver
paths and additive path-weight fold. Its reference meaning is
\[
 D_s(v)=\inf_{p:s\leadsto v}\operatorname{weight}(p)
 \quad\text{in }\mathbb{N}\cup\{\infty\}.
\]
This definition does not invoke the relaxation algorithm. The empty family
of paths has infimum $\infty$; no finite sentinel denotes unreachability.

Initialization assigns $0$ to the source and $\infty$ elsewhere. Relaxing
$e:u\to v$ computes
\[
  c=d(u)+w(e),\qquad d(v)\gets\min(d(v),c).
\]
Every finite label is proved to have an actual source-path witness.
Relaxation preserves that provenance by extending a witnessed path.

Conversely, a store below the initial labels and satisfying every edge
inequality $d(v)\le d(u)+w(e)$ lies below the weight of every path, by path
induction. Therefore a closed store with path provenance equals $D_s$.
Neither queue emptiness without coverage nor a small-looking numerical
assignment is sufficient on its own.

\section{A finite, resumable worklist}

\code{ShortestPathWorklist} retains labels, pending arcs and an ordered
productive-event log. It initially schedules outgoing arcs of the root.
Processing an arc reads the current labels. A strict improvement schedules
the improved vertex's outgoing arcs; a non-improvement consumes only the
current duty.

The invariant has four parts:
\begin{enumerate}
\item labels remain below the initial labels;
\item finite labels retain path provenance;
\item pending arcs belong to the admitted graph;
\item every violated edge inequality has a pending duty.
\end{enumerate}
Only a decrease at a source coordinate can newly violate an outgoing
inequality. That is why rescheduling the updated vertex's outgoing arcs
preserves coverage. At quiescence, coverage proves closure and hence exact
shortest-path distances.

Queue policies preserve duty support, not necessarily duplicate occurrences.
DFS, FIFO and deduplicating FIFO satisfy this condition. Deduplicating pending
work under the final-distance observer is different from deduplicating source
answers. A broken policy that drops newly generated duties terminates with
an incorrect distance and is rejected by the queue law.

Termination uses a lexicographic descent: labels strictly decrease on a
productive step; otherwise the pending-list length decreases. Over finitely
many vertices, the pointwise order on natural distances with top is
well-founded. Hence every finite pending state eventually becomes quiescent.
This is an existence-of-finite-completion theorem, not a uniform numeric
instruction or time bound. The existing \code{MachineCore.runSlice} supplies
the actual bounded runner and exact split-resumption accounting.

\section{The correct lattice direction}

Smaller numerical distances express better upper bounds, so the information
order is the reverse numerical order. \code{LatticeRelaxationClosure}
embeds edge contributions in the existing complete-lattice chaotic-iteration
construction and proves that its least closure equals $D_s$.

Numerically, $D_s$ is the \emph{greatest} feasible label assignment below the
initial labels and satisfying the edge inequalities. Calling it the least
numerical solution of those inequalities would be wrong: unjustified smaller
labels may satisfy them. Path provenance and the information-order reversal
keep the two readings distinct.

Every fair schedule of edge contributions eventually reaches the same
distance assignment. The concrete worklist also handles an empty graph,
where there is no infinite edge-valued schedule to instantiate. No theorem
identifies the schedules' productive-event counts or source round counters.

\section{A stronger law inside one adjacency loop}

The source-shaped \code{relax} function reads a vertex distance once and then
walks its outgoing edges. A live implementation may instead read that source
distance at each edge. \code{RelaxationAdjacency} proves the two loops equal,
including their ordered improvement emissions, for any captured outgoing
edge list with nonnegative weights and no external interference inside the
loop.

An update to a different target leaves the source coordinate alone. A
self-loop cannot improve its own source, because $d(u)+w\ge d(u)$.
Consequently the source value stays fixed throughout this loop. This is
value-level noninterference despite possible syntactic read/write overlap;
requiring every possible write key to differ from every read key would be
unnecessarily restrictive here.

The negative control improves the source externally from $10$ to $0$ before
an edge of weight $1$. The captured and live calculations then propose $11$
and $1$. The isolated-loop law therefore cannot be used to erase an admitted
external update.

\section{Controls, observations and qualification}

On the graph $s\to b$ of cost $1$, $s\to a$ of cost $10$, and $b\to a$ of
cost $1$, DFS and FIFO both finish at $d(b)=1,d(a)=2$. Their chronological
improvement logs are respectively
\[
 [(b,1),(a,2)]\quad\text{and}\quad[(b,1),(a,10),(a,2)].
\]
Thus final-distance equality does not preserve propagation multiplicity.
The inspected source also returns its round counter: a native per-key upsert
can retain the existing driver order, but changing that driver requires a
separate observation theorem or an explicitly different observer contract.

Additional controls cover empty graphs, duplicate edges, unchanged zero-cost
cycles, incomplete prefixes, exact resumption, finite sentinels, stale tickets
and graph deletion invalidating path witnesses. A non-strict improvement guard
requeues a zero-cost self-loop forever. A negative integer self-cycle has no
fixed finite distance. The proof is therefore not a blanket termination
claim for arbitrary lattices, signed graphs, floating arithmetic or PLN
revision. For probabilistic inference, the appropriate information order
must be established; taking the largest scalar truth value is not supplied
by these laws.

The seven mathematical modules and joined audit use the required Lean build
procedure with strict source checks. Module-owned declarations and primitive
axiom dependencies are inventoried; compiler runtime counterparts are listed
separately. Native source extraction, atomic implementation, liveness under
external scheduling and measured costs remain separate qualification duties.
No runtime source is changed by this development.

\end{document}
